\chapter{Contouring}
\label{ch:contour}

\texttt{AMRContour} extracts an isosurface in three dimensions and a contour
curve in two, from the complex rather than from the grid.

\section{Marching simplices, not marching cubes}

The complex is subdivided into simplices---tetrahedra in 3D, triangles in
2D---and contoured with marching tetrahedra (respectively triangles). There are
three distinct cases and no ambiguous faces, so the extracted surface is
topologically consistent by construction. There is no asymptotic-decider
question to answer~\cite{chernyaev1995mc33} and none of the several hundred
lines of case tables that \texttt{isosurface.cpp} carries.

\section{Subdivision: cone from the minimum identifier}
\label{sec:cone}

Each cell is coned from its \emph{minimum-$\gid$} corner. In 3D this means the
three faces not containing the apex, each split on its own minimum-$\gid$
diagonal; in 2D, the two edges not containing it.

Every face of the cell then ends up split by the diagonal through its own
minimum-$\gid$ corner. For the faces containing the apex this is automatic,
since a corner minimal over the cell is minimal over any face containing it;
for the opposite faces it is imposed. Two cells sharing a face see the same
corner identifiers and therefore choose the same diagonal.

\begin{designrule}\label{rule:noparity}
Do not key the diagonal choice on index parity. Parity differs between levels,
so a parity rule fails exactly where two levels meet---the one place the
subdivision must be right.
\end{designrule}

This construction was adopted only after establishing the premise it depends
on. That two cells at a coarse--fine interface share a face \emph{as a matching
set of nodes} is not obvious: the stitch cells' coarse-side faces collapse to
points or edges for most index parities, and a T-junctioned complex would tile
the domain by volume while failing to share faces. The face-incidence census of
\cref{sec:faces} settles it by measurement.

Degenerate simplices---those with repeated identifiers, arising from stitch
cells---are dropped, which is what turns a degenerate hexahedron into the
pyramid or prism it actually is.

\section{Surface vertices are named by their edge}
\label{sec:edgekey}

A surface vertex is identified by the edge it lies on, and that edge is named by
the unordered pair of node identifiers at its ends. Two elements sharing an
edge therefore find the same vertex by identity, at any level and across the
periodic wrap---no coordinate welding, no tolerance, and none of the crack and
mis-placement failures that follow from either.

Everything about a vertex must be a function of the edge alone. This was not
initially true: the position was reconstructed as \emph{origin plus offset} in
the frame of the enclosing element, and two elements sharing an edge have
different origins, so they produced copies differing by an ulp. Within a rank
the cache hid this completely---only the first element to reach an edge ever
computed it---and it surfaced only as two ranks disagreeing. The position is now
computed from the two end nodes' own coordinates, with the far node brought to
its minimum image of the near one:
\begin{equation}
  \mathbf{x}_v \;=\; \mathbf{x}_0 + t\,\bigl(\mathbf{x}_1^{\,\mathrm{mi}} - \mathbf{x}_0\bigr),
  \qquad
  t = \frac{v - f_0}{f_1 - f_0},
  \label{eq:vertpos}
\end{equation}
with the pair ordered by identifier so that $t$ is likewise independent of who
asked.

\begin{designrule}\label{rule:edgefn}
Nothing about the element that reached an edge may enter what its vertex is.
\end{designrule}

\section{Orientation}
\label{sec:orient}

Nothing upstream constrains the winding: vertex order emerges from the marching
case, which follows the simplex's corner order, which follows the cone
construction---all arbitrary relative to the geometry. Measured on a three-level
sphere before this was imposed: \textbf{64598 of 157680 boundary items
misoriented}, about $41\%$, which is what random produces. The corresponding
check in \texttt{isosurface.cpp} is an empty function body.

Each element is now oriented so its normal points toward \emph{increasing} $f$.
The reference direction is from the below-value corners of the simplex to its
above-value corners: exact, local, and independent of the simplex's handedness.

Verification is combinatorial rather than a comparison of normals. A closed
surface is consistently oriented exactly when every interior edge is traversed
in \emph{opposite} directions by the two elements sharing it. That is reported
two ways: a rank-local exact count, and a signed modular sum in which each
traversal contributes $\pm$ its hash, so a correct pair cancels wherever the two
elements live and the global result is exactly zero at any rank count. In 2D the
same statement applies to segment endpoints---each vertex the head of one
segment and the tail of the other.

As an independent geometric confirmation, the enclosed volume by the divergence
theorem on a sphere of radius $0.3$ evaluates to $+0.1130050$ against an exact
$0.1130974$: positive to $0.08\%$, so the normals point outward rather than
merely agreeing with one another.

\section{Carried components}
\label{sec:carried}

Requirement~\ref{req:sample}. Every carried field is interpolated at the same
parameter $t$ of \cref{eq:vertpos} that the isosurface component determines,
which is the convention of \texttt{isosurface.cpp}. Because the vertex is cached
on its edge key, a carried value is computed \emph{once} and every element
sharing that edge receives the identical result---across levels and across the
wrap---rather than each recomputing and the copies being reconciled afterwards
by comparing coordinates.

One caveat is worth recording. At a stitch edge the two endpoints lie on
different levels, so a carried value there is linearly blended across a
resolution jump over a physically longer edge. That is the truthful thing to do
with the data available, and it applies equally to the surface position itself,
but sampled values in the stitch band are a blend across levels rather than a
fine-level quantity.

\section{Output}

MEF is written in the format the existing \texttt{*MEF} chain reads: title,
variable names, element and node-per-element counts, a node block as an
\texttt{FArrayBox} in node-major order, and one-based connectivity. Under MPI it
gathers to the I/O processor---the same wall \texttt{isosurface.cpp} has, and the
reason a \texttt{.pvtu} path remains on the roadmap. It differs in one respect
that matters: duplicate vertices merge by \emph{edge key} rather than by
coordinate proximity, so the merge is exact and the global node ordering, a sort
on those keys, is identical at any rank count.

Both orderings must be canonicalised. Sorting the nodes and not the elements
left the node block byte-identical across rank counts and the connectivity block
differing from its first byte, which is how the defect was localised without
further instrumentation.
